%% kapa: exam=1 round=20 year=2009 period=2 form=A subject=accounting subject_ko=회계학 seq=1 number=1
%% source: Q-Net 감정평가사 기출문제 (cst00309) artlSeq=5024178 file=extracted/1008/5024178/2009년 제20회 감정평가사 1차시험 기출문제/2교시.hwp
%% conversion: kordoc 4.24.51 -> tools/gen1.py (auto); review=proofread-sonnet; flags=none
\begin{question}{1}
다음 중 수익의 인식에 대한 회계처리의 내용으로 옳지 \uline{않은} 것은?
\begin{choices}
  \item 하나의 입장권으로 여러 행사에 참가할 수 있는 경우의 입장료 수익은 각각의 행사를 위해 수행된 용역의 정도에 따라 각 행사에 배분하여 인식한다.
  \item 방송사 등의 광고수익은 해당 광고를 대중에게 전달하는 시점에 수익으로 인식하고, 광고제작사 등의 광고제작용역수익은 진행기준에 따라 수익으로 인식한다.
  \item 거래 이후에 판매자가 소유에 따른 위험을 일부 부담하더라도 그 위험이 별로 중요하지 않은 경우에는 해당 거래를 판매로 보아 수익을 인식한다.
  \item 할부판매의 경우에는 이자수익에 해당하는 부분을 제외한 판매가액을 재화가 인도되는 시점에 수익으로 인식한다.
  \item 상품권의 발행과 관련된 수익은 상품권을 판매한 시점에 수익을 인식하고, 회수하지 못할 것으로 예상되는 부분에 대해서는 충당금을 설정한다.
\end{choices}
\end{question}
\answer{5}
